% titel: Blatt 2, A1: Summenformeln aus Ionen
% skript: Mortimer Kap. 3 + 4 (Salze, Ionen, Formeln)

% Teil I: Werkzeuge, kein Lösungsweg.

Salze sind aus Ionen aufgebaut --- typischerweise aus Kationen (positiv) und
Anionen (negativ). Die Summenformel gibt an, wie viele Ionen jeder Sorte sich zu
einer volkstümlich ladungsneutralen Einheit zusammenfügen.

\begin{werkzeugkasten}[Ladungsneutrale Summenformel]
Ein Salz muss insgesamt ladungsneutral sein: Die Gesamtzahl der positiven
Ladungen muss der Gesamtzahl der negativen Ladungen entsprechen. In der
Bruchteilschreibweise laesst sich das so prüfen: Für ein Salz aus dem Kation
$\mathrm{M}^{z_+}$ und dem Anion $\mathrm{X}^{z_-}$ gilt
\begin{equation}
  z_+ \cdot \nu_- = z_- \cdot \nu_+ \qquad\text{(Ladungsausgleich)},
\end{equation}
wobei $\nu_+$ und $\nu_-$ die kleinsten ganzen Zahlen sind, für die die
Gleichung stimmt. Schlicht: Man braucht so viele Anionen, dass ihre Ladungen die
Ladung des Kations genau ausgleichen, und umgekehrt.

Wichtige Kationen und Anionen (Ladungen aus Namen ableiten):
\begin{itemize}
  \item Eisen-(III) $= \mathrm{Fe^{3+}}$, Aluminium-(III) $= \mathrm{Al^{3+}}$,
        Magnesium $= \mathrm{Mg^{2+}}$, Ammonium $= \mathrm{NH_4^+}$.
  \item Chlorid $= \mathrm{Cl^-}$, Oxid $= \mathrm{O^{2-}}$, Phosphat
        $= \mathrm{PO_4^{3-}}$, Nitrat $= \mathrm{NO_3^-}$, Sulfat
        $= \mathrm{SO_4^{2-}}$; mehratomige Ionen in der Formel in Klammern.
\end{itemize}
\textit{Skript:} Mortimer Kap. 3 (Ionen, Ionengitter) und Kap. 4.1
(Summenformeln), Übung 2/1.
\end{werkzeugkasten}

\begin{achtung}[Darauf achten]
\begin{itemize}
  \item Mehratomige Ionen (z.~B. $\mathrm{PO_4^{3-}}$, $\mathrm{NO_3^-}$,
        $\mathrm{SO_4^{2-}}$) bei Bedarf in Klammern setzen: $\mathrm{Mg(NO_3)_2}$,
        nicht $\mathrm{MgNO_3,2}$.
  \item Die geladenen Ionen selbst, nicht die ungeladenen Elemente einsetzen;
        die Ladung steht im Namen (Eisen-(III) = dreifach positiv).
  \item Kürzen auf das kleinste ganze Zahlenverhaeltnis (nicht 2:4, sondern
        1:2).
  \item Zaehlen beim Aufaddieren: mehratomiges Ion gilt als ein Teilchen.
\end{itemize}
\end{achtung}
